% year: תשע"ט
% type: מבחן
% semester: א
% moed: מיוחד
% number: 1א
% points: 10
% categories: func-limits
% source: exams/מבחנים/תשעט/מועד מיוחד תשעט.pdf

\begin{question}
מצאו את הגבול הבא:
$\displaystyle \lim_{x\to\infty}\left(1+2x-4x^2+x^4\right)^{1/\ln(2x^2+1)}$.
\end{question}

\begin{hint}
כתבו את הביטוי בצורה $e^{\ln(\cdots)/\ln(2x^2+1)}$ והוציאו את החזקה הגבוהה ביותר מתוך כל לוגריתם.
\end{hint}

\begin{solution}
נסמן $p(x)=1+2x-4x^2+x^4$. עבור $x$ גדול מספיק $p(x)>0$ (כי $p(x)\to\infty$), ולכן
\[ p(x)^{1/\ln(2x^2+1)} = \exp\left(\frac{\ln p(x)}{\ln(2x^2+1)}\right). \]
נוציא את החזקה המובילה מכל לוגריתם:
\[ \ln p(x) = \ln\left(x^4\left(1+\tfrac{2}{x^3}-\tfrac{4}{x^2}+\tfrac{1}{x^4}\right)\right) = 4\ln x + \ln\left(1+\tfrac{2}{x^3}-\tfrac{4}{x^2}+\tfrac{1}{x^4}\right), \]
\[ \ln(2x^2+1) = 2\ln x + \ln\left(2+\tfrac{1}{x^2}\right). \]
כאשר $x\to\infty$, הביטויים $\ln\left(1+\tfrac{2}{x^3}-\tfrac{4}{x^2}+\tfrac{1}{x^4}\right)\to \ln 1 = 0$ ו-$\ln\left(2+\tfrac1{x^2}\right)\to\ln 2$ (רציפות הלוגריתם), ואילו $\ln x\to\infty$. נחלק מונה ומכנה ב-$\ln x$:
\[ \frac{\ln p(x)}{\ln(2x^2+1)} = \frac{4 + \frac{\ln\left(1+\frac{2}{x^3}-\frac{4}{x^2}+\frac{1}{x^4}\right)}{\ln x}}{2+\frac{\ln\left(2+\frac{1}{x^2}\right)}{\ln x}} \xrightarrow[x\to\infty]{} \frac{4+0}{2+0}=2. \]
(אפשר גם להשתמש בכלל לופיטל במקרה $\frac{\infty}{\infty}$: $\frac{p'(x)/p(x)}{4x/(2x^2+1)}=\frac{(4x^3-8x+2)(2x^2+1)}{4x\,p(x)}\to\frac{8}{4}=2$.)

מרציפות פונקציית האקספוננט נקבל
\[ \lim_{x\to\infty}\left(1+2x-4x^2+x^4\right)^{1/\ln(2x^2+1)} = e^{2}. \]
\textbf{תשובה:} $e^2$.
\end{solution}
